% ==========================================================
% titles/title-01.tex -- Title I: Foundational Values and Inviolable Rights
% ==========================================================

\ConstTitle{Foundational Values and Inviolable Rights}{title:foundational}

\Article{The Anti-Domination Principle}{art:antidom}

\ArtSection{The Core Prohibition}{art:antidom:s-core-prohibition}
No person, cooperative, Commons institution, government body, algorithm, federation, foreign state, or any other entity may hold arbitrary power over the survival or self-determination of any person. This principle is the interpretive foundation of all other provisions within this Constitution.

\ArtSection{Ownership and Domination}{art:antidom:s-ownership-domination}
Ownership is legitimate when it reflects direct use and does not create a system of power over the survival of others. Ownership that creates dependency relationships -- compelling others to submit to arbitrary control in order to meet their basic needs -- is not a protected right under this Constitution.

\ArtSection{The Scope of This Principle}{art:antidom:s-scope-principle}
The anti-domination principle applies equally to:
\begin{itemize}
  \item Individual relationships between persons
  \item Economic relationships between worker-owners, cooperatives, and federations
  \item Institutional relationships between governance bodies and citizens
  \item International relationships between this Commonwealth and other nations and peoples
\end{itemize}

\Article{The Material Sufficiency Floor}{art:matsuff}

\ArtSection{Guaranteed Material Rights}{art:matsuff:s-guaranteed-material-rights}
The following are inviolable rights of every person present on Commonwealth territory, unconditionally and without means-testing:
\begin{itemize}
  \item Adequate housing, maintained to a standard of safety, habitability, and dignity
  \item Food stability sufficient for health and wellbeing
  \item Comprehensive healthcare, including physical, mental, reproductive, gender-affirming, and preventive care
  \item Education at all levels from early childhood through advanced study
  \item Childcare of genuine quality, accessible to all who require it
  \item Elder care and care for persons with disabilities
\end{itemize}

\ArtSection{The Basis of This Guarantee}{art:matsuff:s-basis-guarantee}
These rights exist because their absence places persons in conditions of material desperation that make domination possible. A person who may be deprived of food, shelter, or healthcare can be controlled through the threat of that deprivation. This Constitution aims to remove that lever entirely.

\ArtSection{Universality}{art:matsuff:s-universality}
These rights apply to every person on Commonwealth territory regardless of citizenship status, employment status, criminal record, disability, age, or any other characteristic. They are not earned, means-tested, or revocable. They are the material precondition of freedom. These rights apply without diminishment to persons in secure facilities. Incarceration does not suspend, reduce, or modify any right guaranteed in this Article. See also Article \ref{art:justice-restoration}, \ref{art:justice-restoration:s-secure-facilities}.

\ArtSection{Immutability}{art:matsuff:s-immutability}
The Material Sufficiency Floor may be expanded but never reduced. No constitutional revision, emergency declaration, legislative act, judicial decision, or executive action may diminish the rights guaranteed in this Article. This is one of three truly hard limits in this Constitution. See Article \ref{art:const-assembly}, \ref{art:hardlimits}.

\Article{Equality and Bodily Autonomy}{art:equality}

\ArtSection{Racial Equality}{art:equality:s-racial-equality}
Race is a morally arbitrary characteristic that has been systematically used as a mechanism of domination and exploitation. All persons are equal regardless of race, ethnicity, or national origin. The Commonwealth acknowledges its historical debts to communities harmed by racial domination and commits to active reparative justice, the terms of which are determined through the processes established in Article \ref{art:reparations}.

\ArtSection{Gender Equality}{art:equality:s-gender-equality}
All persons are equal regardless of sex assigned at birth. The Commonwealth recognizes that reproductive and care labor is the foundation upon which all other labor rests. Its socialization through the Material Sufficiency Floor is not welfare policy but the correction of a structural extraction that prior economic systems imposed disproportionately on individuals who were assigned Female at birth.

\ArtSection{Sexual Orientation and Gender Identity}{art:equality:s-sexual-orientation-gender}
All persons are equal regardless of sexual orientation or gender identity. The state has no legitimate interest in the intimate life, relational choices, or gender identity of any consenting adult. All legal recognitions -- partnership, marriage, adoption, inheritance, and all others -- are available to any consenting adults without restriction based on gender or orientation. Gender-affirming care is part of the Healthcare Commons.

\ArtSection{Reproductive Autonomy}{art:equality:s-reproductive-autonomy}
Every person possesses inviolable autonomy over their own body and reproduction. The state may not compel pregnancy, prohibit abortion, mandate sterilization, or otherwise substitute its judgment for the individual's on questions of reproduction. The full spectrum of reproductive healthcare -- contraception, abortion, fertility treatment, prenatal and postnatal care -- is included within the Healthcare Commons without moral conditions imposed by providers or administrators.

\ArtSection{Other Protected Characteristics}{art:equality:s-other-protected-characteristics}
No person may face domination, discrimination, or diminished access to any right or Commons service on the basis of disability, age, religion, language, or any other characteristic not relevant to the specific matter at hand.

\ArtSection{The Distinction Between Belief and Imposition}{art:equality:s-distinction-between-belief}
Freedom of belief, conscience, and internal community practice is absolute. No person or community may be compelled to abandon their convictions. However, no belief entitles any person or community to use state power, economic leverage, or communal governance to control the identity, body, relationships, or reproduction of others. The boundary between internal practice and external imposition is enforced by the \gls{gls:constitutional-court}.

\ArtSection{Separation of Governance and Religion}{art:equality:s-separation-governance-religion}
The Commonwealth is a secular institution. No governance body, official document, currency, public oath, or official communication of the Commonwealth invokes, endorses, references, or incorporates any religious belief, practice, symbol, or language. This prohibition applies to all layers of governance without exception. Freedom of religious belief and practice by individuals and communities is protected under Section 5 of this Article; this Section governs the Commonwealth's own conduct, not the private conduct of its members.

\ArtSection{Freedom from Compelled Expression}{art:equality:s-freedom-compelled-expression}
No person may be required by any Commonwealth institution, educational body, employer, or governance body to recite, display, perform, or affirm any pledge, oath, anthem, or symbol of national loyalty as a condition of participation, membership, education, employment, or any other benefit or right. Voluntary participation in expressions of civic culture is protected; compelled participation is prohibited. This prohibition applies to all educational settings without exception, from early childhood through advanced study.

\Article{The Anti-Entrenchment Principle}{art:antientrench}

\ArtSection{The Core Prohibition}{art:antientrench:s-core-prohibition}
No person, party, institution, cooperative, Commons body, algorithm, or governance structure may make itself permanent, unreviewable, or insulated from democratic revision. All authority in this Commonwealth is temporary, delegated, and revocable.

\ArtSection{Application}{art:antientrench:s-application}
The anti-entrenchment principle applies to all layers of governance, including the \gls{gls:constitutional-assembly} itself. No assembly, legislature, court, or executive body may take actions whose primary effect is to make its own continuation more likely or its own decisions harder to reverse.

\ArtSection{Dissolution Prohibition}{art:antientrench:s-dissolution-prohibition}
No person or group of persons, regardless of the offices they hold, may dissolve or fundamentally restructure the constitutional order without the full processes established in Articles \ref{art:ordinary-amend} and \ref{art:failure}. Any purported dissolution or restructuring outside those processes is void. Upon any such attempt, the \gls{gls:transparency-engine} immediately issues a public alert to all residents, the \gls{gls:constitutional-court} automatically convenes within twenty-four hours, and a mandatory national referendum is automatically scheduled within sixty days.

\Article{Radical Transparency}{art:transparency}

\ArtSection{The Principle}{art:transparency:s-principle}
All exercises of power at every level of governance -- deliberations, decisions, communications, resource flows, and reasoning -- are made official, and publicly accessible in real time, through the \gls{gls:transparency-engine}. Secrecy in governance is the exception and requires specific constitutional authorization; transparency is the default and requires no justification.

\ArtSection{Limits on Transparency}{art:transparency:s-limits-transparency}
Personal privacy is protected. Operational security matters may be temporarily withheld subject to constitutional authorization but may not be withheld for more than five years. All withheld matters must be logged as withheld in the public record even while their content is protected. Military capability classifications under the Capability Classification Protocol of Article \ref{art:defense}, \ref{art:defense:s-defensive-posture-review} are subject to the same five-year maximum withholding limit.